\chapter{Point Defects}

No crystal is perfect. Point defects---vacancies, interstitials, and impurity atoms---are present in every real crystal and have profound effects on mechanical, electrical, and optical properties.

\section{Lattice Vacancies}

A vacancy is an empty lattice site. The equilibrium concentration at temperature $T$ is:
\begin{equation}
  n_v = N\exp\left(-\frac{E_v}{k_BT}\right)
\end{equation}
where $E_v$ is the vacancy formation energy.

\begin{table}[H]
\centering
\caption{Vacancy formation energies in metals, from experimental measurements (positron annihilation, differential thermal expansion, quenching experiments).}
\label{tab:vacancy}
\begin{tabular}{@{}lSS@{}}
\toprule
Metal & {$E_v$ (\si{\electronvolt})} & {$n_v/N$ at melting point} \\
\midrule
Cu & 1.28 & 2e-4 \\
Ag & 1.09 & 2e-4 \\
Au & 0.94 & 7e-4 \\
Al & 0.75 & 9e-4 \\
Pb & 0.49 & 2e-3 \\
Na & 0.42 & 7e-3 \\
Fe & 1.60 & 1e-4 \\
W  & 3.60 & 4e-7 \\
\bottomrule
\end{tabular}
\end{table}

The vacancy concentration at the melting point ranges from $\sim 10^{-3}$ (Na, Pb---low $E_v$) to $\sim 10^{-7}$ (W---very high $E_v$). Even the highest concentrations correspond to only one vacancy per $\sim$1000 atoms.

\section{Color Centers}

\subsection{F Centers}

An \textbf{F center} (from the German \textit{Farbe}, color) is an electron trapped at an anion vacancy in an alkali halide crystal. The trapped electron absorbs visible light, giving the crystal a characteristic color.

Seitz (1946)~\cite{seitz1946fcenters} wrote the definitive review of color centers (612 citations), establishing the F center as a model quantum system---essentially a particle in a box formed by the surrounding cations.

\begin{table}[H]
\centering
\caption{F-center absorption peak wavelengths and energies in alkali halides.}
\label{tab:fcenter}
\begin{tabular}{@{}lSSl@{}}
\toprule
Crystal & {$\lambda_{\max}$ (\si{\nano\metre})} & {$E$ (\si{\electronvolt})} & Color \\
\midrule
LiF  & 250 & 4.96 & colorless (UV) \\
NaF  & 340 & 3.65 & colorless (UV) \\
NaCl & 465 & 2.67 & yellow-brown \\
NaBr & 540 & 2.30 & brown \\
KCl  & 563 & 2.20 & violet \\
KBr  & 620 & 2.00 & blue-green \\
KI   & 690 & 1.80 & green \\
RbCl & 620 & 2.00 & blue \\
\bottomrule
\end{tabular}
\end{table}

The F-center absorption shifts to longer wavelengths (lower energies) as the lattice parameter increases---the \textbf{Mollwo--Ivey relation}: $E \propto a^{-n}$ with $n \approx 1.8$. This is a direct consequence of the larger ``box'' (anion vacancy) in crystals with bigger lattice constants.

\vspace{1cm}
\noindent\rule{\textwidth}{0.4pt}
\textit{Extended line defects---dislocations---determine the mechanical properties of real crystals.}
